事件函数把应用的一段处理逻辑交给平台按事件执行。
应用指定入口函数 handler，提交代码及依赖，配置执行角色和资源，
再把 HTTP 入口、队列或其他事件来源接到函数上。
平台负责准备执行环境、启动运行时、把事件交给 handler，
并根据调用方式处理响应、排队和重新交付。
这种由服务商统一管理执行服务器、容量和维护的方式，
通常称为 Serverless~\cite{lambda-event-overview}。

AWS Lambda 提供这种托管计算服务。
购买系统可以把事件函数用作查询入口、事件消费者或恢复工作者。
一次执行读取原业务身份，访问所属权威与外部平台，再提交本次处理结果。
函数返回、调用方获知、外部效果确定和权威保存原结果，
分别发生在这条链路的不同位置。

从系统设计看，部署配置决定运行什么代码、使用哪些权限与资源；
事件接入决定谁交付事件、谁等待结果；
执行环境承载运行时、handler 和局部缓存；
业务权威与平台账本保存跨执行环境的责任和效果。
图~\ref{fig:ch20-lambda-platform} 按这些公开职责表示事件函数路径。
部署单元是代码与配置，扩容单元是执行环境，
长期恢复单元则由原 $C/T$、$EV$ 或其他持久业务身份确定。

\CTwentyFigure{ch20-32}{Lambda 事件函数的部署、接入、执行环境与业务保存域}{fig:ch20-lambda-platform}

本节按取得的官方说明，展开 Event Functions
在默认容量提供者下的普通事件执行路径。
同一服务还提供以 HTTP 服务器处理请求的 Web Functions，
以及为一个用户或作业提供独立环境的 MicroVMs 等形态，
它们的调用、并发与寿命有各自规则。
Event Functions 中的 Durable Functions
增加可持久恢复的多步骤工作流，
Managed Instances 则使用另外的实例容量与并发安排。
以下的同步、异步和事件源映射，
按所讨论的事件函数路径说明~\cite{lambda-event-overview,lambda-lifecycle}。

\subsection{同步调用把响应交回调用者}

同步调用中，调用者向 Lambda 发送事件，
服务安排函数执行并把函数响应交回调用者。
当调用 API 成功运行了函数时，
API 可以返回 HTTP 200，
同时通过函数错误字段和响应内容表达处理错误。
调用者分别检查服务请求是否成功以及函数报告了什么结果~\cite{lambda-invocation}。

调用方可能在函数结束前超时或断开。
它没有取得结果时，函数可能尚未开始、正在执行，
或者已经完成平台付款。
决定重试后，事件继续携带原 $S$ 和原输入，
handler 读取原权威绑定，
返回原结果或接续原 $C/T$。
SDK 对超时、限流或服务错误的重试，也通过这一业务身份关联。

同步函数错误之后由谁再次提交调用，
取决于调用方及其上游配置。
API Gateway、其他服务或应用 SDK 各自拥有等待和重试规则，
异步事件队列的重试安排属于另一条路径。
部署图画出真正发起下一次请求的主体，
就能确定由谁保存原输入、控制等待并重新交付事件。

\CTwentyFigure{ch20-16}{同步、异步与事件源映射分别拥有不同的重试主体}{fig:ch20-lambda-invocation}

\subsection{异步调用先进入服务队列}

异步调用先把事件放入 Lambda 的异步队列，
接收响应表示事件已被接受，函数执行随后发生。
函数错误默认再尝试两次，
第一次错误后通常等待约一分钟，
第二次错误后再等待约两分钟。
限流和系统错误采用另外的退避及事件年龄规则，
配置可以调整重试次数、最长保留时间和失败去向~\cite{lambda-invocation}。

即使函数没有报告错误，同一事件也可能重复交付。
队列压力还可能让事件在真正执行前达到年龄上限。
死信队列或失败目标保存后续处理入口，
恢复者从中取得原业务身份，再查询当前权威。
失败去向里的一条记录描述相应交付或执行遇到的问题，
原 $C$ 若已成功，恢复者从权威或平台读回成功结果。

设函数第一次已按 \code{pay-C1} 形成 $X_1$，
返回过程失败，异步队列随后再交付事件。
第二次 handler 从权威或平台原操作取得 $X_1$，
补齐原结果后返回。
如果处理程序以本次尝试编号创建付款键，
正常的服务重试就会形成第二个平台操作。
原键从持久责任传入 handler，
才使不同执行尝试共同访问原来的收款事实。

\subsection{事件源映射按记录与批次推进}

事件源映射是 Lambda 管理的一类资源。
其中的 poller 即轮询执行者，从流或队列读取记录，
再按批次调用函数。
S3、SNS 等服务直接触发函数的路径，
与这种主动拉取来源数据的方式各有接口。
批次大小、批次窗口和负载大小共同决定何时形成调用，
不同来源的默认值与推进规则也不同~\cite{lambda-event-source}。

流来源涉及分片位置与顺序。
失败批次可能阻止该分片的后续处理，
重试、拆分与丢弃由来源和配置决定。
Amazon SQS 是托管消息队列，
围绕消息暂时不可见、成功删除和超时后重新可见组织处理。
读取 SQS 的轮询者与读取流分片的轮询者，
分别面对可见性超时和消费位置等进度事实。

设一个 SQS 批次包含消息 a、b、c，
a 和 c 已在消费事务中保存 $I/F$，b 失败。
整批返回错误时，a、c 也可能再次可见。
它们重新到达后，消费域读取原 Inbox 与 business guard，
返回原效果关联。
这两个记录分别保存消息接收与业务效果约束，
让重复交付从已提交的消费事务继续。

\CTwentyFigure{ch20-17}{整批错误与部分批次响应采用不同的重新交付范围}{fig:ch20-lambda-batch}

支持的 SQS 配置启用 \code{ReportBatchItemFailures} 后，
handler 可以在 \code{batchItemFailures} 中只列失败消息标识。
上例返回 b 的消息 ID，
服务据此让失败项重新进入后续处理，
成功项按正常完成路径处理。
若函数直接抛出异常，整个批次被视为失败。
FIFO 队列要保持先进先出的相应处理顺序，
遇到首个失败后停止继续处理，
并将失败项及尚未处理项一并报告~\cite{lambda-sqs-partial}。

部分批次响应减少不必要的重投。
消费事务继续保存完整 $I$、原 payload、业务 guard、
$F$ 与 dispatch 意图。
c 若已提交消费事务、函数尚未返回部分失败结果就退出，
服务可能再次交付 c，消费域从原记录回答。
消息 ID 关联来源记录，批次关联一次组合交付，
函数调用 ID 关联执行，业务效果身份关联原履约责任。
它们分别保存各层进度，恢复时再沿记录关联取得完整事实。

\subsection{部署包怎样进入执行环境}

前面的重投会把同一事件交给新的执行尝试。
部署与环境结构决定这些尝试能共享哪些内容：
代码版本可以相同，进程内存和可写目录却随环境分开。
先看代码怎样装入环境，再看环境复用、接口响应和隔离边界，
就能确定一次中断后哪些局部状态可能留下，哪些材料要重新读取。

事件函数可以使用 ZIP 部署包，也可以使用容器镜像。
部署包带入函数代码与依赖，
运行时负责加载代码并调用 handler。
容器镜像将运行时、依赖和入口组织成可部署的文件系统内容，
平台再为函数建立执行环境。
镜像版本、某个环境实例和一次事件调用具有不同身份，
更新镜像形成新的代码内容，
扩容则可能让多个环境同时运行同一版本~\cite{lambda-container-images}。

Lambda 容器镜像必须实现 Runtime API 所要求的运行时接口，
可以使用官方提供的语言运行时与接口客户端，
也可以把接口客户端加入自选的 Linux 基础镜像。
镜像上传到同地域的 Amazon ECR 镜像仓库后，
函数配置引用相应镜像。
镜像面向单一支持的处理器架构，
代码和依赖要能在只读文件系统下运行，
可写临时内容放在环境的 \code{/tmp} 中。
部署者据此把代码版本、系统依赖和可写缓存分别放置。

\CTwentyFigure{ch20-33}{镜像承载代码与依赖，执行环境承载进程和局部可写状态}{fig:ch20-lambda-package}

CPU、内存、执行超时和并发配置共同约束处理能力。
在这里的默认事件函数路径中，一个环境同一时段承载一次调用，
新增并发调用可能要求建立更多环境。
保留并发为某函数划出账户并发份额并限制其使用上限，
预置并发则提前准备已初始化环境。
两者分别控制可取得的执行数量和初始化等待，
支付平台及权威数据库的下游容量另由它们自己的配额和准入控制~\cite{lambda-concurrency}。

\subsection{冷启动、复用与重置}

执行环境的 Init 阶段建立扩展、运行时与函数静态代码，
Invoke 阶段运行 handler，
Shutdown 阶段结束相应进程。
没有可复用环境而需要建立新环境时，
初始化工作进入冷启动路径。
复用已有环境可以减少这部分开销，
连接池和只读数据适合放在恰当的初始化位置~\cite{lambda-lifecycle}。

同一环境可能连续处理多次调用，
全局变量和部分临时文件会被后续调用看到。
环境也可能被回收，
执行失败后可能重置运行时，再重新初始化。
SnapStart 通过保存初始化后的快照并恢复环境，
预置并发通过预先准备环境减少等待，
它们与普通按需初始化具有不同的启动步骤。
使用 Managed Instances 等方式时，
再根据对应并发和生命周期规则安排局部状态。

\CTwentyFigure{ch20-18}{执行环境可复用也可重置，持久业务责任跨越环境生命周期}{fig:ch20-lambda-lifecycle}

全局字典记录已经处理过 $C_1$ 时，
环境替换会使这份内存消失，
多个环境之间也各有自己的字典。
反过来，复用环境若错误保留上一次账户的变量，
下一次调用又可能读到不属于本次输入的内容。
缓存按完整身份和版本组织，
受理绑定与原结果从持久权威取得，
每次 handler 据此建立本次处理上下文。

\code{/tmp} 具有执行环境局部性，
某些重置过程还会保留其中内容。
新环境则未必取得旧环境文件，
所以事件指向持久受理记录，
或携带可核验的原身份关联，
handler 再从权威读取原 \code{scope/pay-C} 和当前责任。
快照恢复后的连接、随机状态与身份分配，
也按对应运行时规则重新处理。
原 $C/T$ 已在业务权威分配时，恢复环境直接读取它们。

\subsection{Runtime API 怎样连接运行时}

Runtime API 是 Lambda 环境内服务与语言运行时之间的接口。
自定义运行时从 \code{AWS\_LAMBDA\_RUNTIME\_API}
取得地址，调用
\code{GET /2018-06-01/runtime/invocation/next}
等待下一事件。
响应体提供事件，响应头提供 request ID、
截止时间、函数 ARN 与跟踪信息。
这里的函数 ARN 标识被调用的函数、版本或别名。
运行时将事件和上下文交给 handler，
再把结果 POST 到对应 request ID 的 response 路径，
错误则提交到相应 error 路径~\cite{lambda-runtime-api}。

request ID 关联触发调用的事件，
一个 request ID 可以对应多次调用尝试。
invocation ID 则标识其中一次尝试。
运行时接到这个头后，在 response 或 error 中回传，
服务据此核对当前尝试。
官方接口为兼容已有运行时允许省略该头，
传入与当前尝试不匹配的值则进入相应拒绝分支。
这些字段使运行时响应归入正确的事件与尝试，
原付款结果则通过业务 $C$ 和平台键取得。

\CTwentyFigure{ch20-19}{Runtime API 关联事件、当前调用尝试与返回路径}{fig:ch20-runtime-api}

运行时等待 \code{next}、handler 访问权威、
handler 请求平台以及运行时提交 response，
经过不同接口和保存域。
handler 得到付款结果后、运行时提交 response 前退出，
平台账本可以已保存 $X_1$，
权威可能已回填，也可能尚未取得确定结果。
新环境按原 $C/T$ 核对两边事实，
再返回原结果或继续查询。

\CTwentyFigure{ch20-34}{平台提交、权威回填和运行时响应分别发生，超时保留原责任}{fig:ch20-function-commits}

Runtime Interface Client，简称 RIC，
负责相应语言的运行时循环、事件处理及 Runtime API 交互。
Runtime Interface Emulator，简称 RIE，
在本地提供相关接口代理，
让容器化函数通过 HTTP 接收测试事件。
Python RIC 4.2.2 与 RIE 1.37 有公开实现。
RIE 面向 Linux，
其本地接口验证与云端编排、安全、认证属于不同范围~\cite{lambda-ric-rie}。

handler 的输入与结果可以在普通进程中检查，
Runtime API 协议可以验证事件接收和响应关联，
RIC/RIE 则承载相应官方本地运行路径。
云端事件重试、账户配额、环境回收、权限与隔离，
由实际服务配置及部署形成。

\subsection{Firecracker 承担的隔离边界}

函数平台要在共享物理机上运行来自不同工作负载的代码，
同时控制单个工作负载能访问的内存、设备与资源。
容器通过宿主内核提供的进程和资源隔离组织程序，
虚拟机则为客体提供独立内核及虚拟硬件边界。
短函数启动频繁、单次工作较少，
执行平台希望在保留虚拟机边界的同时，降低启动与常驻开销。

Firecracker 是面向这类容器与函数负载的轻量虚拟机监控器，
即 VMM。
它使用宿主 Linux 的 KVM 硬件虚拟化基础设施建立 MicroVM，
每个 MicroVM 对应一个 Firecracker 进程，
内部运行客体内核与用户程序。
VMM 保留较小的设备集合，
用 virtio 为客体提供网络和块设备 I/O。
virtio 通过客体与宿主约定的虚拟设备接口传递请求，
让客体程序访问受配置约束的网络与存储~\cite{firecracker-design}。

部署者通过配置 API 指定客体内存、CPU、内核、块设备与网络参数，
启动后由客体加载运行时和代码。
jailer 是约束 VMM 进程的外层组件，
通过受限文件系统、命名空间、降权和系统调用过滤等方式收紧宿主访问。
宿主资源控制与设备 I/O 限速，
还限制单个工作负载占用共享机器的程度。
虚拟机边界、VMM 外层约束与函数调用接口由此分别对应具体组件。

\CTwentyFigure{ch20-20}{MicroVM、VMM 与宿主约束执行边界，业务事实保存在独立权威}{fig:ch20-firecracker}

Firecracker 的公开部署设计给出了 Lambda 事件路径的一种完整组织：
接入前端进行认证、授权并取得函数元数据，
路由管理者选择已有执行槽，
没有可用槽时由放置服务选择 worker 建立新槽。
执行槽表示预先加载了某函数运行环境的执行位置，
同一槽可以在先后多次调用中复用。
worker 上的 MicroVM 为这个执行位置提供隔离，
放置与路由协作则组织资源寿命、负载和重复利用。
图~\ref{fig:ch20-slot-design} 按这份公开设计画出角色关系\footnote{这里采用
Firecracker 集中书目所列公开设计中的事件路径和 worker 结构，
属于其披露的 Lambda 架构。当前服务的其他容量和执行方式，
按对应官方接口判断。}。

\CTwentyFigure{ch20-35}{公开执行槽设计将接入、路由、放置和 worker 复用连接起来}{fig:ch20-slot-design}

公开 VMM 的代码说明怎样建立与约束一台 MicroVM，
当前 Lambda 接口说明应用可以取得什么服务能力，
两者分别提供组件实现与托管服务的观察入口。
自行部署 Firecracker 时，
部署者还要安排代码分发、路由、权限和故障恢复。
使用托管事件函数时，这些执行安排由平台承载，
业务持久责任则通过应用权威和外部账本保持关联。

MicroVM重新建立后，客体内存与本地状态按启动或快照恢复路径形成。
假定快照形成在handler发送付款以前，而平台随后已保存$X_1$，
从这份快照启动的环境就同时面对旧局部状态与较新的平台账本。
它从持久任务取得原\code{scope/pay-C}，先查询该操作，
便能把恢复位置与平台已经发生的收款重新对齐。
VMM的内存和设备边界说明这份局部状态在哪里，
恢复协议说明它怎样接回环境之外的业务历史。

\subsection{短执行单元怎样承担长期恢复}

一次函数调用开始时，从持久任务 $T$ 读取原 $C$、
原 \code{scope/pay-C} 与发送资格。
先核对联合库存 $TX$ 的持久 commit 条件，
再取得当前执行资格，
查询或发送允许的原平台操作。
确定结果回写原 $C/T$，
与原结果、$O/EV$ 在所属权威提交中关联。
执行接近时限而结果尚未确定时，
权威继续保存未决责任，后继调用再取得这份工作。

如果一次调用只能处理一批待办，
扫描游标的持久推进要与每项处置对应。
已经结束的 $C/T$ 保存终态，
未决项继续留在可恢复集合，
游标表示哪些记录已经完成本轮归类。
后继调用从重叠区间扫描，
已完成项返回原结果，未决项按资格继续。
这样，减少重复扫描的进度信息与长期责任集合分别承担自己的作用。

\CTwentyFigure{ch20-21}{函数调用结束后，未决 $C/T$ 和原平台身份继续约束下一次恢复}{fig:ch20-function-continuation}

事件函数、Ray 任务、Actor 与批量作业可以使用同一业务恢复入口，
执行成本则各不相同。
大量小任务增加调度和序列化开销，
过大的 handler 批次扩大超时后的重复处理范围，
过高恢复并发会挤占在线权威与平台配额。
部署时分别预算新受理、同 $S$ 重试、外部尝试与未决扫描的资源，
为每类工作安排队列、并发上限和实际执行机会。

\let\code\BeforeTwentyCode
\setlength{\emergencystretch}{\BeforeTwentyEmergency}
\tolerance=\BeforeTwentyTolerance
